\section{The master shorted--Gram gate}
\label{sec:tpc62-master-shorted}

The shorted-operator construction and its Loewner-maximal
characterization are classical, originating with Anderson and the
Anderson--Trapp development
\citep{Anderson1971,AndersonTrapp1975}.  We include the finite-dimensional
proof to fix the exact quotient, nuisance space, and quantifiers needed by
the TPC closure specialization; no priority claim is made for the general
shorted-operator theory.

This section isolates the finite-dimensional closure problem before any
arithmetic specialization.  Let \(\cK\) and \(\cH\) be finite-dimensional
complex Hilbert spaces, let
\begin{equation}
 R=R^*\ge 0\quad\text{on }\cK,
 \qquad
 F:\cK\longrightarrow\cH,
 \label{eq:tpc62-master-data}
\end{equation}
and assume \(R\ne0\).  The parent form is
\(z\mapsto\langle Rz,z\rangle\), while \(F\) is the final physical map
after every declared layer has been assembled.  Define the exact closure
constant
\begin{equation}
 \alpha_R(F)
 :=\inf_{\langle Rz,z\rangle>0}
 \frac{\|Fz\|^2}{\langle Rz,z\rangle}.
 \label{eq:tpc62-alpha-definition}
\end{equation}
We also write \(\alpha(F;R):=\alpha_R(F)\) when the map is the first
argument in a chain of certificates.
The order of the quantifiers in \eqref{eq:tpc62-alpha-definition} is part
of the definition: one constant must control every parent-visible vector.
The exclusion \(R=0\) is substantive rather than cosmetic.  In that case
there is no parent-visible vector, the infimum in
\eqref{eq:tpc62-alpha-definition} has empty domain, and the least
eigenvalue on \(S=\{0\}\) is not defined.  Every inequality
\(\|Fz\|^2\ge c\langle Rz,z\rangle\) is then vacuous, so the zero parent
form is treated separately and is assigned no closure constant.

Put
\begin{equation}
 N:=\Ker R,
 \qquad S:=N^\perp=\Ran R,
 \qquad R_S:=R|_S,
 \label{eq:tpc62-master-NSR}
\end{equation}
so that \(R_S\) is positive definite on the nonzero space \(S\).  Write
\begin{equation}
 F_S:=F|_S,
 \qquad F_N:=F|_N,
 \qquad W:=F(N)=\Ran F_N,
 \qquad P:=P_{W^\perp}.
 \label{eq:tpc62-master-FSNWP}
\end{equation}
The subspace \(W\) is important: vectors in \(N\) cost no parent energy,
but their final images may cancel the image of a parent-visible vector.

\subsection{The finite-dimensional short}

For a positive semidefinite operator \(G\) on \(\cK\) and a subspace
\(L\subseteq\cK\), its \emph{short to \(L\)}, denoted \(\Sh_L(G)\), is
the unique positive operator characterized by
\begin{equation}
 \Sh_L(G)
 =\max\bigl\{B:0\le B\le G,\ \Ran B\subseteq L\bigr\}
 \label{eq:tpc62-short-order-definition}
\end{equation}
in the Loewner order.  For completeness, it has the following exact
variational description.  If \(s\in L\), then
\begin{equation}
 \left\langle\Sh_L(G)s,s\right\rangle
 =\min_{n\in L^\perp}\langle G(s+n),s+n\rangle,
 \label{eq:tpc62-short-variational}
\end{equation}
and \(\Sh_L(G)\) vanishes on \(L^\perp\).

To verify both existence and maximality without invoking an
infinite-dimensional theorem, write \(G=C^*C\) and put
\(V=C(L^\perp)\).  Least squares gives
\begin{equation}
 \min_{n\in L^\perp}\|Cs+Cn\|^2
 =\|P_{V^\perp}Cs\|^2.
 \label{eq:tpc62-short-least-squares}
\end{equation}
The right-hand side is a positive quadratic form on \(L\), and therefore
defines a positive operator supported on \(L\).  Call it \(B_L\).  For
\(z=s+n\), \eqref{eq:tpc62-short-variational} gives
\(\langle B_Lz,z\rangle\le\langle Gz,z\rangle\), so \(B_L\le G\).
Conversely, if \(0\le B\le G\) and \(\Ran B\subseteq L\), then
\(B=P_LBP_L\), and for every \(n\in L^\perp\),
\[
 \langle Bs,s\rangle
 =\langle B(s+n),s+n\rangle
 \le\langle G(s+n),s+n\rangle.
\]
Taking the minimum over \(n\) shows \(B\le B_L\).  Hence
\(B_L=\Sh_L(G)\), proving \eqref{eq:tpc62-short-order-definition} and
\eqref{eq:tpc62-short-variational} simultaneously.

Apply this construction to the final Gram operator
\begin{equation}
 G:=F^*F.
 \label{eq:tpc62-final-Gram}
\end{equation}
Relative to \(\cK=S\oplus N\), write
\begin{equation}
 G=
 \begin{pmatrix}
  G_{SS}&G_{SN}\\
  G_{NS}&G_{NN}
 \end{pmatrix}
 =
 \begin{pmatrix}
  F_S^*F_S&F_S^*F_N\\
  F_N^*F_S&F_N^*F_N
 \end{pmatrix}.
 \label{eq:tpc62-Gram-blocks}
\end{equation}

\begin{lemma}[Projected Gram and generalized Schur complement]
\label{lem:tpc62-projected-Gram}
The short of \(G\) to \(S\) is
\begin{equation}
 \Sh_S(G)
 =\begin{pmatrix}Q&0\\0&0\end{pmatrix},
 \qquad
 \boxed{Q=F_S^*PF_S.}
 \label{eq:tpc62-master-Q-projected}
\end{equation}
Equivalently,
\begin{equation}
 Q
 =G_{SS}-G_{SN}G_{NN}^{\dagger}G_{NS}.
 \label{eq:tpc62-master-Q-Schur}
\end{equation}
For each \(s\in S\),
\begin{equation}
 \min_{n\in N}\|F_Ss+F_Nn\|^2
 =\|PF_Ss\|^2
 =\langle Qs,s\rangle.
 \label{eq:tpc62-master-nuisance-minimum}
\end{equation}
All minimizers are
\begin{equation}
 n_s=-F_N^\dagger F_Ss+n_0,
 \qquad n_0\in\Ker F_N.
 \label{eq:tpc62-master-minimizers}
\end{equation}
\end{lemma}

\begin{proof}
The first equality in \eqref{eq:tpc62-master-nuisance-minimum} is the
least-squares distance from \(-F_Ss\) to
\(W=\Ran F_N\).  The Moore--Penrose least-squares formula gives
\eqref{eq:tpc62-master-minimizers}.  Consequently the quadratic form of
the short is \(\|PF_Ss\|^2\), proving
\eqref{eq:tpc62-master-Q-projected}.  Finally,
\[
 F_N(F_N^*F_N)^\dagger F_N^*=P_W,
\]
and substitution in the block expression gives
\eqref{eq:tpc62-master-Q-Schur}.
\end{proof}

\subsection{The master eigenvalue and exact kernel decision}

\begin{theorem}[Master shorted--Gram formula]
\label{thm:tpc62-master-shorted-gram}
For the data in \eqref{eq:tpc62-master-data},
\begin{equation}
 \boxed{
 \alpha_R(F)
 =\lambda_{\min}\!\left(
 R_S^{-1/2}QR_S^{-1/2}\bigm|_S
 \right),
 \qquad W=F(\Ker R),\quad Q=F_S^*P_{W^\perp}F_S.}
 \label{eq:tpc62-master-eigenvalue}
\end{equation}
Equivalently,
\begin{equation}
 \alpha_R(F)
 =\lambda_{\min}\!\left(
 \left.
 R^{\dagger/2}\Sh_S(F^*F)R^{\dagger/2}
 \right|_S
 \right).
 \label{eq:tpc62-master-global-eigenvalue}
\end{equation}
Moreover,
\begin{equation}
 \boxed{
 \alpha_R(F)>0
 \quad\Longleftrightarrow\quad
 \Ker F\subseteq\Ker R.}
 \label{eq:tpc62-master-kernel-equivalence}
\end{equation}
Thus failure of the kernel inclusion is an exact finite-scale kill, not a
loss in a sufficient estimate.
\end{theorem}

\begin{proof}
Every \(z\in\cK\) has a unique decomposition \(z=s+n\), with
\(s\in S\) and \(n\in N\).  Since \(Rn=0\),
\[
 \langle Rz,z\rangle=\langle R_Ss,s\rangle,
\]
which is positive precisely when \(s\ne0\).  Minimizing the numerator over
\(n\) and using \cref{lem:tpc62-projected-Gram} gives
\begin{equation}
 \alpha_R(F)
 =\inf_{s\in S\setminus\{0\}}
 \frac{\langle Qs,s\rangle}{\langle R_Ss,s\rangle}.
 \label{eq:tpc62-master-Rayleigh-quotient}
\end{equation}
The substitution \(x=R_S^{1/2}s\) and the Rayleigh--Ritz principle prove
\eqref{eq:tpc62-master-eigenvalue}.  On \(S\),
\(R^{\dagger/2}=R_S^{-1/2}\), while both
\(R^{\dagger/2}\) and \(\Sh_S(F^*F)\) vanish on \(N\); this proves
\eqref{eq:tpc62-master-global-eigenvalue}.

If \(\Ker F\not\subseteq\Ker R\), choose \(z\in\Ker F\) with
\(\langle Rz,z\rangle>0\).  Its quotient is zero, so
\(\alpha_R(F)=0\).  Conversely, suppose
\(\Ker F\subseteq N\).  If \(PF_Ss=0\), then
\(F_Ss\in W\), so there is \(n\in N\) with
\(F_Ss+F_Nn=0\).  Hence \(s+n\in\Ker F\subseteq N\).  Projecting this
vector orthogonally onto \(S\) gives \(s=0\).  Thus \(PF_S\) is
injective on \(S\).  In finite dimension its least generalized singular
value relative to the positive definite form \(R_S\) is strictly
positive.  Equation \eqref{eq:tpc62-master-eigenvalue} now gives
\(\alpha_R(F)>0\).
\end{proof}

\begin{corollary}[Dimension kill]
\label{cor:tpc62-master-dimension-kill}
If
\begin{equation}
 \dim S>\dim\cH-\dim W,
 \label{eq:tpc62-master-dimension-kill}
\end{equation}
then \(\alpha_R(F)=0\).
\end{corollary}

\begin{proof}
The effective map \(PF_S:S\to W^\perp\) cannot be injective under
\eqref{eq:tpc62-master-dimension-kill}; apply
\eqref{eq:tpc62-master-kernel-equivalence}.
\end{proof}

The converse to \cref{cor:tpc62-master-dimension-kill} is false: compatible
dimensions do not force injectivity.  Even the exact kernel inclusion has
no asymptotic size content.  For example, with \(R_X=I\) and
\(F_X=e^{-X}I\), one has \(\Ker F_X=\{0\}\) and
\begin{equation}
 \alpha_{R_X}(F_X)=e^{-2X}>0
 \label{eq:tpc62-positive-superpolynomial-small}
\end{equation}
at every finite scale, but there is no finite exponent \(\lambda\) for
which \(\alpha_{R_X}(F_X)\ge X^{-\lambda+o(1)}\).  Kernel survival and an
endpoint-quality polynomial certificate are therefore separate gates.

\subsection{A mandatory factor: an actual kill versus a failed certificate}

Suppose the final map has a genuine factorization
\begin{equation}
 \cK\xrightarrow{\ T\ }\mathcal Y\xrightarrow{\ L\ }\cH,
 \qquad F=LT.
 \label{eq:tpc62-mandatory-factorization}
\end{equation}
Calling \(T\) \emph{mandatory} means the identity
\eqref{eq:tpc62-mandatory-factorization} holds for the literal final map;
it does not mean merely that one proof happens to use \(T\).

\begin{proposition}[Mandatory-factor kernel kill]
\label{prop:tpc62-mandatory-factor-kill}
If
\begin{equation}
 \Ker T\not\subseteq\Ker R,
 \label{eq:tpc62-T-visible-kernel}
\end{equation}
then
\begin{equation}
 \alpha_R(LT)=0
 \label{eq:tpc62-T-kills-alpha}
\end{equation}
for every downstream map \(L\).  This conclusion requires an actual
vector in \(\Ker T\setminus\Ker R\); failure to prove a positive bound for
\(T\), or failure of one factorwise sufficient certificate, is not such a
vector.
\end{proposition}

\begin{proof}
Choose \(z\in\Ker T\setminus\Ker R\).  Positivity of \(R\) implies
\(\langle Rz,z\rangle>0\), whereas \(LTz=0\).  Thus the quotient in
\eqref{eq:tpc62-alpha-definition} vanishes.  Equivalently,
\(\Ker T\subseteq\Ker(LT)\), so the kernel criterion
\eqref{eq:tpc62-master-kernel-equivalence} fires.
\end{proof}

The converse of \cref{prop:tpc62-mandatory-factor-kill} is deliberately
not asserted.  The inclusion \(\Ker T\subseteq\Ker R\) says only that the
mandatory factor itself has not killed a parent-visible vector.  A
downstream map can still do so: for \(R=I\), \(T=I\), and \(L=0\), the
inclusion holds while \(\alpha_R(LT)=0\).

Here is a precise sufficient certificate whose failure illustrates the
logical distinction.  When \(T\ne0\), define
\begin{align}
 \tau_R(T)
 &:=\inf_{\langle Rz,z\rangle>0}
   \frac{\|Tz\|^2}{\langle Rz,z\rangle},
 \label{eq:tpc62-tau-T}\\
 \ell(L;T)
 &:=\inf_{0\ne y\in\Ran T}
   \frac{\|Ly\|^2}{\|y\|^2}.
 \label{eq:tpc62-ell-LT}
\end{align}
Then, directly from the two definitions,
\begin{equation}
 \alpha_R(LT)\ge \ell(L;T)\,\tau_R(T).
 \label{eq:tpc62-factorwise-certificate}
\end{equation}
The equality \(\tau_R(T)=0\) is an exact mandatory-factor kill precisely
when \(\Ker T\not\subseteq\Ker R\), by
\cref{thm:tpc62-master-shorted-gram} applied to \(T\).  In contrast,
\(\ell(L;T)=0\) may only say that \(L\) has a kernel on a direction which
the parent form already regards as invisible.

\begin{example}[A factorwise certificate can fail while closure is exact]
\label{ex:tpc62-factor-certificate-failure}
Take \(\cK=\mathcal Y=\C^2\), \(\cH=\C\), and
\begin{equation}
 R=\begin{pmatrix}1&0\\0&0\end{pmatrix},
 \qquad T=I,
 \qquad L(x,y)=x.
 \label{eq:tpc62-factor-certificate-example}
\end{equation}
Here \(\tau_R(T)=1\), but \(\ell(L;T)=0\) because
\(L(0,1)=0\).  Thus the right-hand side of
\eqref{eq:tpc62-factorwise-certificate} is zero.  Nevertheless
\begin{equation}
 \frac{\|LT(x,y)\|^2}{\langle R(x,y),(x,y)\rangle}=1
 \qquad(x\ne0),
 \label{eq:tpc62-factor-certificate-example-ratio}
\end{equation}
so \(\alpha_R(LT)=1\).  Indeed
\(\Ker(LT)=\operatorname{span}\{(0,1)\}=\Ker R\).
The certificate failed; the gate did not.
\end{example}

At an asymptotic scale, the same distinction must be maintained.  An
exhibited vector in \(\Ker T_X\setminus\Ker R_X\) is a rigorous zero at
that scale.  A nonpositive exponent budget, an uncontrolled norm, or the
absence of a proof of \(\tau_{R_X}(T_X)\gg X^{-C}\) is only a failure of
the chosen route.  It cannot be promoted to an operator kernel statement.
